\documentclass[10pt,a4paper]{amsart}
\usepackage[margin=19mm]{geometry}
\usepackage{amsmath,amssymb,mathtools}
\usepackage[scr=rsfs,cal=euler]{mathalfa}
\usepackage{xfrac}
\usepackage[unicode,hidelinks,bookmarks=false]{hyperref}
\newcommand{\ie}{i.e.\ }
\newcommand{\cf}{cf.\ }
\newcommand{\resp}{respectively\ }
\newcommand{\Subsection}{\subsection}
\providecommand{\nobreakdash}{\nobreak\hbox{-}}
\setlength{\emergencystretch}{2em}
\multlinegap0pt
\usepackage{bm}
\usepackage{bbm}
\usepackage{dsfont}
\usepackage{xspace}
\usepackage{cancel}
\usepackage{mathtools}
\usepackage{amsthm}
\makeatletter
\@ifundefined{theorem}{\newtheorem{theorem}{Theorem}}{}
\@ifundefined{proposition}{\newtheorem{proposition}[theorem]{Proposition}}{}
\makeatother
\makeatletter
\expandafter\ifx\csname Proof\endcsname\relax
\newenvironment{Proof}
% {\noindent{\it Proof\/}}{{ \hfill $\Box$}\smallskip\par}

\fi
\expandafter\ifx\csname definition\endcsname\relax
\newenvironment{definition}
% {\smallskip\noindent{\bf Definition\/}:}{\smallskip\par}
% \newenvironment{statement}
% {\smallskip\noindent{\bf Statement\/}.}{\smallskip\par}
% \newenvironment{example}
% {\smallskip\noindent{\bf Example\/}.}{\smallskip\par}
% \newenvironment{remark}
% {\smallskip\noindent{\bf Remark\/}.}{\smallskip\par}
% \newenvironment{remarks}
% {\smallskip\noindent{\bf Remarks\/}.}{\smallskip\par}
% \newenvironment{proof}
% {\noindent{\it Proof\/}.}{{ \hfill $\Box$}\smallskip\par}
% \newenvironment{Proof}
% {\noindent{\it Proof\/}}{{ \hfill $\Box$}\smallskip\par}

\fi
\expandafter\ifx\csname example\endcsname\relax
\newenvironment{example}
% {\smallskip\noindent{\bf Example\/}.}{\smallskip\par}
% \newenvironment{remark}
% {\smallskip\noindent{\bf Remark\/}.}{\smallskip\par}
% \newenvironment{remarks}
% {\smallskip\noindent{\bf Remarks\/}.}{\smallskip\par}
% \newenvironment{proof}
% {\noindent{\it Proof\/}.}{{ \hfill $\Box$}\smallskip\par}
% \newenvironment{Proof}
% {\noindent{\it Proof\/}}{{ \hfill $\Box$}\smallskip\par}

\fi
\expandafter\ifx\csname namelist\endcsname\relax
\newenvironment{namelist}[1]{%
\begin{list}{}
{\let\makelabel\namelistlabel
\settowidth{\labelwidth}{#1}
\setlength{\leftmargin}{1.1\labelwidth}}
}{%
\end{list}}
\fi
\expandafter\ifx\csname proof\endcsname\relax
\newenvironment{proof}
% {\noindent{\it Proof\/}.}{{ \hfill $\Box$}\smallskip\par}
% \newenvironment{Proof}
% {\noindent{\it Proof\/}}{{ \hfill $\Box$}\smallskip\par}

\fi
\expandafter\ifx\csname remark\endcsname\relax
\newenvironment{remark}
% {\smallskip\noindent{\bf Remark\/}.}{\smallskip\par}
% \newenvironment{remarks}
% {\smallskip\noindent{\bf Remarks\/}.}{\smallskip\par}
% \newenvironment{proof}
% {\noindent{\it Proof\/}.}{{ \hfill $\Box$}\smallskip\par}
% \newenvironment{Proof}
% {\noindent{\it Proof\/}}{{ \hfill $\Box$}\smallskip\par}

\fi
\expandafter\ifx\csname remarks\endcsname\relax
\newenvironment{remarks}
% {\smallskip\noindent{\bf Remarks\/}.}{\smallskip\par}
% \newenvironment{proof}
% {\noindent{\it Proof\/}.}{{ \hfill $\Box$}\smallskip\par}
% \newenvironment{Proof}
% {\noindent{\it Proof\/}}{{ \hfill $\Box$}\smallskip\par}

\fi
\expandafter\ifx\csname statement\endcsname\relax
\newenvironment{statement}
% {\smallskip\noindent{\bf Statement\/}.}{\smallskip\par}
% \newenvironment{example}
% {\smallskip\noindent{\bf Example\/}.}{\smallskip\par}
% \newenvironment{remark}
% {\smallskip\noindent{\bf Remark\/}.}{\smallskip\par}
% \newenvironment{remarks}
% {\smallskip\noindent{\bf Remarks\/}.}{\smallskip\par}
% \newenvironment{proof}
% {\noindent{\it Proof\/}.}{{ \hfill $\Box$}\smallskip\par}
% \newenvironment{Proof}
% {\noindent{\it Proof\/}}{{ \hfill $\Box$}\smallskip\par}

\fi
\makeatother
\title[Non-invertible quasihomogeneous singularities and their Land]{Non-invertible quasihomogeneous singularities and their Landau-Ginzburg orbifolds}
\author{Anton Rarovskii}
\date{Source: arXiv:2405.17091v2 (CC BY 4.0)}
\begin{document}
\maketitle

\noindent\emph{Source and licence.} Non-invertible quasihomogeneous singularities and their Landau-Ginzburg orbifolds. Version arXiv:2405.17091v2 (CC BY 4.0), distributed under the Creative Commons Attribution 4.0 International licence, as recorded in the arXiv licence metadata for this exact version and shown on the abstract page https://arxiv.org/abs/2405.17091v2. Full rights notice: licenses/arxiv-2405.17091-CC-BY-4.0.txt.\par\smallskip

\noindent\textit{Editorial scope.} Counted unit: the Proposition whose source label is \texttt{None} in the article (source0.tex lines 925--927, source label \texttt{None}), delivered together with the article's own complete original proof of it (source0.tex lines 928--966, one \texttt{proof} environment of 39 lines). No mathematics is written, repaired or re-derived: statement and proof are the article's own text line for line. Required context retained uncounted: none. Every definition, display and label of those lines is retained unchanged. Omitted from this package and not redistributed: 1--924, 967--1167. Nothing inside a retained excerpt is omitted or altered: every retained body is delivered exactly as the article's lines are written, so no omission receipt of any kind accompanies this package. Citations: the retained excerpts carry no citation command at all, so no bibliography entry of the article is transcribed and no reference is redistributed. Reference adaptation: none was needed and none was made; every reference inside the delivered unit resolves inside it, so no unresolved reference or citation is produced. Printed slips kept literal: no printed word, symbol, hypothesis or number of the counted statement or of its proof was altered. Layout and class: the article's own class is \texttt{amsart} with packages including amsmath, amsthm, amssymb, amsfonts, amscd, epsfig, xy, tabularx, booktabs, float, mathtext, fontenc, multicol, quiver; the wrapper is the lane's pinned \texttt{amsart} editorial wrapper at 10pt on A4, which restates the theorem-like environments the retained text needs and adds the article's own macro definitions that the retained text needs (none), with extra wrapper packages (bm, bbm, dsfont, xspace, cancel, mathtools). The delivered numbering is the unit's own. Assets: no retained span contains a figure environment, a graphics-inclusion command, a PDF-page inclusion, a table or a TikZ picture; no image file is copied into this package, the package declares no asset dependency and no image file is redistributed with it. Licence: version arXiv:2405.17091v2 of this article is distributed under the Creative Commons Attribution 4.0 International licence, as recorded in the arXiv licence metadata for this exact version and shown on the abstract page https://arxiv.org/abs/2405.17091v2. A full rights notice is delivered at licenses/arxiv-2405.17091-CC-BY-4.0.txt. Characters: every retained excerpt is pure ASCII, so the delivered engine, XeLaTeX with its default Latin Modern text font, reports no missing character.
\begin{proposition}
     All critical points of $\hat{f}$ are on the chart $U_{1}$. Moreover, $\hat{f}_{1}$ is non-degenerate.
\end{proposition}

\begin{proof}
Our aim is to prove that $0$ is an isolated solution of the system $\{d \hat{f}_{s} = 0\}$ only if $s=1$. Recall the explicit form of $\hat{f}_{1}$ and put $y = x_1$ and $z=x_{N+1}$:
\begin{align*}
  \hat{f}_{1} = (\hat{f}_\kappa)_{1} + (\hat{f}_{add})_{1} = f_{\kappa_0}(x_1,\dots,x_N) + x_{N+1}^2x_1 +& \sum_{\{J_r \in A_R\;|\;
1 \notin J_r\}} \varepsilon_{J_r}x_{s_1}^{b_{s_1}}x_{s_2}^{b_{s_2}} \dots x_{s_l}^{b_{s_l}} +\\ +& \sum_{\{J_r \in A_R\;|\;
1 \in J_r\}} \varepsilon_{J_r}x_{1}^{b_1}x_{s_2}^{b_{s_2}} \dots x_{s_l}^{b_{s_l}}
\end{align*}
It is not hard to note that if $s$ such that $1 < s < N+1$ and $s$ does not lie in any set from $A_R$, then
\[\frac{\partial \hat{f}_{1}}{\partial x_s} =  \frac{\partial f_{\kappa_0}}{\partial x_s} \]
So in this case $0$ is a solution of equations above since $f$ is non-degenerate.
If $s$ such that $1 < s < N+1$ and $s\in J_k$ for some $J_k \in A_R$, then the system of equations
\[\frac{\partial \hat{f}_1}{\partial x_s} = a_sx_s^{a_s-1}t_{\kappa(s)} + \sum_{w \in \kappa^{-1}(s)}x_w^{a_w} + \sum_{J_k \in A_R \wedge s \in J_k} \varepsilon_{J_k}b_{s}x_{s}^{b_{s}-1}x_{k_2}^{b_{k_2}} \dots x_{k_l}^{b_{k_l}} \]
obviously vanishes at $0$. And for $s = 1,N+1$ we also have $0$ as a solution:
\[
\frac{\partial \hat{f}_{1}}{\partial x_1} = a_1x_1^{a_1-1}x_{2} + x_{N+1}^2 + \sum_{\{J_r \in A_R\;|\;
1 \in J_r\}} \varepsilon_{J_r}2b_1x_{1}^{2b_1-1}x_{s_2}^{b_{s_2}} \dots x_{s_l}^{b_{s_l}} = 0
\]
\[
\frac{\partial \hat{f}_1}{\partial x_{N+1}} = x_1x_{N+1} = 0
\]

We have shown that $0$ is a solution of the equations above. Also note that the change of variables $\{t_1 = x_1, t_{N+1}^2 = x_1x_{N+1}^2, t_1t_{N+1}= x_1x_{N+1}\}$ is a diffeomorphism and since $f$ is non-degenerate then  $\hat{f}_1$ is also non-degenerate.

Let us turn to $\hat{f}_2$, rewrite it similarly to $\hat{f}_1$ and consider two equations:
\[
\frac{\partial \hat{f}_2}{\partial x_1} = 2a_1x_1^{2a_1-1}x_{N+1}^{a_{N+1}}x_{2} + \sum_{J_k \in A_R \wedge 1 \in J_k} \varepsilon_{J_k}2b_{k_1}x_1^{2b_{k_1}-1}x_{N+1}^{b_{k_1}}x_{k_2}^{b_{k_2}} \dots x_{k_l}^{b_{k_l}} = 0
\]
\[
\frac{\partial \hat{f}_2}{\partial x_{N+1}} = a_ix_1^{2a_1}x_{N+1}^{a_1-1}x_{2} + \sum_{J_k \in A_R \wedge 1 \in J_k} \varepsilon_{J_k}b_{k_1}x_1^{2b_{k_1}}x_{N+1}^{b_{k_1}-1}x_{k_2}^{b_{k_2}} \dots x_{k_l}^{b_{k_l}} + 1 = 0
\]
Obviously this equations does not have the solution at the origin. Moreover, LHS of the second equation does not vanish if $x_1=0$ or $x_{N+1}=0$. It means we can suppose that $x_1\neq 0$ and $x_{N+1} \neq 0$. Then we divide the first equation on $2x_1^{2a_1-1}x_{N+1}^{a_1}$ and the second one on $x_1^{2a_1}x_{N+1}^{a_1-1}$, and obtain the following:
\[
a_1x_{2} + \sum_{J_k \in A_R \wedge 1 \in J_k} \varepsilon_{J_k}b_{k_1}x_1^{2b_{k_1}-2a_1}x_{N+1}^{b_{k_1}-a_1}x_{k_2}^{b_{k_2}} \dots x_{k_l}^{b_{k_l}} = 0
\]
\[
a_1x_{2} + \sum_{J_k \in A_R \wedge 1 \in J_k} \varepsilon_{J_k} b_{k_1}x_1^{2b_{k_1}-2a_1}x_{N+1}^{b_{k_1}-a_1}x_{k_2}^{b_{k_2}} \dots x_{k_l}^{b_{k_l}} + x_1^{-2a_1}x_{N+1}^{1-a_1} = 0
\]
which implies $x_1^{-2a_1}x_{N+1}^{1-a_1} = 0$. Thus, the system $\{d\hat{f}_2 = 0\}$ does not have any solution what completes the proof.
\end{proof}

\end{document}
